
% -----------------------------------------------------------------
% ---------------------- P Ř Í K L A D Y --------------------------
% -----------------------------------------------------------------

\exercisepart
\setcounter{section}{1}

% =========================
% ÚLOHA 1 – Ověřte správnost tvrzení o logaritmech
% =========================
\begin{exercise}[Ověřte správnost následujících tvrzení]
    Pomocí definice logaritmu ověřte, zda jsou následující tvrzení pravdivá:
    \begin{tasks}(2)
        \task $\log_6 6 = 1$
        \task $\log_7 49 = 2$
        \task $\log 100 = 2$
        \task $\log_2 4 = 16$
        \task $\log_2 0{,}5 = -1$
        \task $\log_9 \sqrt{3} = 0{,}25$
        \task $\log_{27} \sqrt{27} = 0{,}125$
        \task $\log_5 25 = 3$
        \task $\log 0{,}01 = -2$
        \task $\log_3 81 = 4$
        \task $\log_4 16 = 3$
        \task $\log_{0{,}1} \frac{1}{10} = -1$
    \end{tasks}
\end{exercise}


% =========================
% ÚLOHA 2 – Vypočítejte hodnoty logaritmů
% =========================
\begin{exercise}[Vypočítejte hodnoty následujících logaritmů]
    Určete hodnotu každého z následujících logaritmů. Použijte definici dekadického logaritmu
    a základní vzorce pro práci s logaritmy.
    \begin{tasks}(2)
        \task $\log 10$
        \task $\log 100$
        \task $\log 10^{5}$
        \task $\log \sqrt{10}$
        \task $\log 10^{-3}$
        \task $\log 1$
        \task $\log 0{,}1$
        \task $\log 0{,}01$
    \end{tasks}
\end{exercise}


% =========================
% ÚLOHA 3 – Vypočítejte hodnoty logaritmických výrazů
% =========================
\begin{exercise}[Vypočítejte hodnoty následujících výrazů]
    Určete hodnotu výrazů využívajících základní ekvivalenci.
    \begin{tasks}(2)
        \task $5^{\log_5 2}$
        \task $10^{\log_{10} 1}$
        \task $8^{\,1+\log_{8} 5}$
        \task $3^{2\log_{3} 4}$
        \task $2^{\log_{2} 0{,}25}$
        \task $7^{\log_{7} 14 - \log_{7} 2}$
        \task $5^{\log_{5} 20 - \log_{5} 4}$
        \task $4^{\log_{4} 50 + \log_{4} 2}$
        \task $9^{\log_{9} 3^{3}}$
        \task $6^{\log_{6} 0{,}1}$
        \task $10^{\log 5 + \log 2}$
        \task $2^{\log_2 32 - \log_2 4}$
    \end{tasks}
\end{exercise}


% =========================
% ÚLOHA 4 – Jednoduché rovnice s logaritmem
% =========================
\begin{exercise}[Najděte všechny hodnoty proměnné]
    Najděte všechny hodnoty $x \in (0; +\infty)$, pro něž platí:
    \begin{tasks}(2)
        \task $\log_2 x = 0$
        \task $\log_2 x = 1$
        \task $\log_2 x = -3$
        \task $\log_2 x = -\frac{1}{5}$
        \task $\log_5 x = -3$
        \task $\log_5 x = 1$
        \task $\log_5 x = 0$
        \task $\log_5 x = -\frac{1}{3}$
        \task $\log_{10} x = 4$
        \task $\log_{10} x = -2$
        \task $\log_{10} x = 0$
        \task $\log_{10} x = 1$
    \end{tasks}
\end{exercise}


% =========================
% ÚLOHA 5 – Hledání vhodného základu logaritmu
% =========================
\begin{exercise}[Určete všechny hodnoty parametru $a$]
    Určete všechna reálná čísla $a$, pro něž platí:
    \begin{tasks}(2)
        \task $\log_a 4 = 2$
        \task $\log_a 0{,}0001 = -2$
        \task $\log_a 25 = 2$
        \task $\log_a 0{,}01 = -2$
        \task $\log_a 16 = 4$
        \task $\log_a 0{,}0001 = -4$
        \task $\log_a 9 = 2$
        \task $\log_a 0{,}09 = -2$
        \task $\log_a 81 = 4$
        \task $\log_a 0{,}0001 = -4$
    \end{tasks}
\end{exercise}


% =========================
% ÚLOHA 6 – Bonus: Náčrtky grafů netradičních funkcí
% =========================
\begin{exercise}[BONUS: Načrtněte grafy následujících funkcí]
    Načrtněte grafy těchto funkcí (uveďte také jejich definiční obory):
    \begin{tasks}(2)
        \task $y = \log_x 2$
        \task $y = 2^{\log_2 x}$
        \task $y = \log_x x$
        \task $y = x^{\log_x x}$
        \task $y = \log_x \left(\dfrac{1}{x}\right)$
        \task $y = \log_x \left(x^{\sqrt{x}}\right)$
        \task $y = \log_x \left(\log_x x\right)$
        \task $\star \; y = \log_x x \cdot \bigl(\log_x x^{\log_x (x^{\sqrt{x}})}\bigr)$
    \end{tasks}
\end{exercise}


% =========================
% ÚLOHA 7 – Upravte logaritmické výrazy
% =========================
\begin{exercise}[Upravte a vypočítejte hodnotu výrazů]
    Využijte vzorce pro součin, podíl a mocninu uvnitř logaritmu.
    \begin{tasks}(2)
        \task $4\log_6 3 \;+\; 5\log_6 2 \;-\; \log_6 12$
        \task $\log 1500 \;-\; \log 15$
        \task $\log_5 10 \;+\; \log_5 12{,}5$
        \task $-\log_3 4 \;+\; \log_3 6 \;-\; \dfrac{1}{4} \log_3 0{,}5$
        \task $\dfrac{2}{5} \left( 4\log_6 11 \;-\; \dfrac{1}{6} \log_6 4 \;+\; \log_6 5 \right)$
        \task $3\log_2 5 \;+\; 2\log_2 3 \;-\; \log_2 60$
        \task $\log_4 32 \;-\; 2\log_4 3 \;+\; \dfrac{1}{2} \log_4 9$
        \task $5\log_7 2 \;-\; 3\log_7 4 \;+\; \log_7 14$
        \task $2\log_3 5 \;+\; \log_3 9 \;-\; \log_3 20$
        \task $\log_2 40 \;-\; \log_2 5 \;+\; \dfrac{1}{2} \log_2 16$
        \task $-\log_8 27 \;+\; \log_8 9 \;+\; 2\log_8 6$
        \task $\dfrac{3}{2} \log_4 8 \;-\; \log_4 2 \;+\; \log_4 32$
    \end{tasks}
\end{exercise}


% =========================
% ÚLOHA 8 – Exponenciální pokles koncentrace DDT
% =========================
\begin{exercise}[Rozklad škodlivé látky DDT v potravním řetězci]

    \begin{minipage}[t]{0.68\textwidth}
        Látka DDT (dichlordifenyltrichlorethan), pro člověka velmi škodlivá, se hromadí 
        v potravním řetězci a dostává se tak do mléka a dalších potravin.
        Její koncentrace ve výši $5 \cdot 10^{-6}\,\%$ je v současné době ještě tolerovaná,
        do budoucna je však požadováno snížit ji na $2 \cdot 10^{-6}\,\%$.
        Používání DDT je dnes téměř ve všech státech zakázáno, ale chemický rozklad probíhá
        jen velmi pozvolna — poločas rozkladu je přibližně $30$ let.
    \end{minipage}
    \hfill
    \begin{minipage}[t]{0.28\textwidth}
        \begin{center}
          \smallskip
            \qrcode[height=2cm]{https://www.instagram.com/postrehy_z_wikipedie/p/C-aZd6siaIe/}\\[4pt]
            {\small Příspěvek k tématu na Instagramu, proč se zakázalo používat DDT.}
        \end{center}
    \end{minipage}

    \medskip

    \begin{tasks}(1)
        \task Předpokládejme, že pokles koncentrace lze popsat vztahem
        \[
            C(t) = C_0 \left(\frac{1}{2}\right)^{\frac{t}{30}},
        \]
        kde $C_0$ je počáteční koncentrace a $t$ čas v letech.  
        Zapište rovnici, kterou musí splňovat čas $t$, kdy koncentrace klesne z 
        $5 \cdot 10^{-6}\,\%$ na $2 \cdot 10^{-6}\,\%$.

        \task Vyřešte tuto rovnici a určete, za jak dlouho (na celé roky) 
        bude dosaženo požadované nižší koncentrace.
    \end{tasks}
\end{exercise}


% =========================
% ÚLOHA 9 – Složitější logaritmické rovnice
% =========================
\begin{exercise}[Řešte složitější logaritmické rovnice]
    Řešte následující rovnice, určete všechny reálné kořeny.
    \begin{tasks}(1)
        \task $\log_{2}(x+1) = 3$
        \task $4\log_{3}(2x-1) = 12$
        \task $\log_{5}(x-2) + \log_{5}(x+3) = 2$
        \task $2\log_{4}(x+1) - \log_{4}(x-3) = 2$
        \task $\log_{\frac12}(2-x) = -2$
        \task $\log_{4}(5x-4) = 2$
        \task $\log_{2}\log_{3}\log_{\frac12} x = 0$
        \task $\log_{\frac12}\log_{3}\bigl(1 + 20\log_{2} x\bigr) = -2$
        \task $\log_{2}\bigl[14 + 2\log_{7}\bigl(1 + 2\log_{\frac12} x\bigr)\bigr] = 4$
        \task $\log_{9}\Bigl\{3\log_{2}\bigl[1 + \log_{3}\bigl(1 - 2\log_{3} x\bigr)\bigr]\Bigr\} = 0{,}5$
        \task $\log_{3}\bigl[5 + \log_{2}(8 - x)\bigr] = 2$
        \task $\log_{\frac12}\Bigl\{4 + \log_{5}\bigl[10 + \log_{2}(x - 3)\bigr]\Bigr\} = -3$
        \task $\log_{5}\Bigl\{6 + \log_{2}\bigl[8 + \log_{3}(x - 1)\bigr]\Bigr\} = 2$
    \end{tasks}
\end{exercise}


% =========================
% ÚLOHA 10 – Řešte logaritmické rovnice
% =========================
\begin{exercise}[Řešte následující logaritmické rovnice]
    Určete všechna reálná čísla $x$, pro která mají rovnice smysl a která dané rovnice splňují.
    \begin{tasks}(2)
        \task $\log_3(x+5) = \log_3(2x-1)$
        \task $\log_5(x^2 - 17) = \log_5(x+3)$
        \task $\log_2(3x - 4) + \log_2(x + 1) = 3$
        \task $2\log_4(x - 2) - \log_4(x + 6) = 1$
        \task $\log_{0{,}5}(x + 1) + \log_{0{,}5}(x - 3) = 2$
        \task $3\log_2(x - 1) + \log_2(x + 3) = 5$
        \task $\log_7(x^2 - 4) - \log_7(x - 2) = 1$
        \task $2\log_3(x + 4) - \log_3(x - 1) = 2$
        \task $\log_{2}(x - 3) + \log_{2}(x + 5) = 4$
        \task $ \log_{5}(2x + 1) - \log_{5}(x - 1) = 1$
        \task $\log_{0{,}2}(x + 2) + \log_{0{,}2}(x - 4) = -3$
        \task $ \log_{4}(x^2 - 9) - \log_{4}(x + 3) = 1$
    \end{tasks}
\end{exercise}


% =========================
% ÚLOHA 11 – Další logaritmické rovnice
% =========================
\begin{exercise}[Řešte logaritmické rovnice s neznámou $x \in \mathbb{R}$]
    Najděte všechna reálná čísla $x$, pro která mají rovnice smysl a která je splňují.
    \begin{tasks}(2)
        \task $\log x = 2 - \log 5$
        \task $\log_{2}(x+14) + \log_{2}(x+2) = 6$
        \task $\log x = 2\log 5 + \log 4$
        \task $\log_{5}(x-1) + \log_{5}(x+3) = 2$
        \task $3\log_{2}(x+1) - \log_{2}(x-3) = 2$
        \task $\log_{3}(x+4) + \log_{9}(x-2) = 3$  
        \task $\log_{10}(x+3) - \log_{10} 2 = 1 - \log_{10}(x+2)$
        \task $\log_{4}(2x+5) - \log_{4}(x+1) = 1$
        \task $\log_{6}(x+1) + \log_{3} x = 1$
        \task $\log_{8}\sqrt{x+30} + \log_{8}\sqrt{x} = 1$
        \task \text{$\log_{10}\sqrt{x-4} + \log_{10}\sqrt{3x+1} = \log_{10} 40 - 1$}
        \task $\dfrac{2\log_{10}(x+2)}{\log_{10}(3x+10)} = 1$
        \task \text{$\log\sqrt{x} + \log\!\left(\dfrac{1}{x^{2}}\right) - \log x^{3} + \dfrac{11}{2}
               = \dfrac{\log x^{2}}{1 + \log 10}$}
    \end{tasks}
\end{exercise}


% =========================
% ÚLOHA 12 – Logaritmické rovnice s racionálními výrazy
% =========================
\begin{exercise}[Řešte logaritmické rovnice]
    Určete všechna reálná čísla $x$, pro která mají rovnice smysl a která dané rovnice splňují.

    \begin{tasks}(2)
        \task $\log_{2}\!\left(\frac{3 - x}{x + 3}\right) = -2$
        \task $\log_{5}\!\left(\frac{2x + 1}{x - 4}\right) = 1$
        \task $\log_{10}\!\left(\frac{x - 2}{x + 8}\right) = -1$
        \task $\log_{3}\!\left(\frac{6x - 2}{x - 3}\right) = 2$
        \task $\log_{\frac12}\!\left(\frac{x}{x + 14}\right) = \dfrac{\log 125}{\log 5}$
        \task $\log_{7}\!\left(\frac{x^{2} + 1}{x^{2} - 1}\right) = 1$
        \task $\log_{4}\!\left(\frac{x-2}{x+6}\right) + \log_{4}(x-1) = 1$
        \task $2\log_{3}(x+4) - \log_{3}(x^{2}-1) = \log_{3} 3$
        \task $\log_{0{,}5}\!\left(\frac{x+5}{2x-1}\right) - \log_{0{,}5}(x-3) = -2$
        \task $\log_{2}\!\left(\log_{3}(x+1)\right) = \log_{2}\!\left(\dfrac{1}{\,\log_{3}(x-2)}\right)$
        \task $\log_{5}\!\left(\frac{x^{2} - 4x + 4}{x^{2} - 9}\right) = 0$
        \task $ \log_{10}\!\left(\frac{2x + 6}{x - 3}\right) + \log_{10}(x - 5) = 1$
        \task $\log_{3}\!\left(\frac{x + 1}{x - 2}\right) + \log_{3}(x - 1) = 2$
        \task $\log_{2}\!\left(\frac{x^{2} - 5x + 6}{x^{2} - 4}\right) = 3$
    \end{tasks}
\end{exercise}


% =========================
% ÚLOHA 13 – Rovnice s logaritmy v zlomku
% =========================
\begin{exercise}[Řešte rovnice s neznámou]
    Řešte v množině reálných čísel $x$:
    \begin{tasks}(2)
        \task $\displaystyle \frac{\log_3(6x-2)}{\log_3(x-3)} = 2$
        \task $\displaystyle \frac{\log_5\left(x - \frac{1}{4}\right)}{\log_5\left(x + \frac{7}{2}\right)} = -1$
        \task $\displaystyle \frac{2\log(3x)}{\log(2-7x)} = 1$
        \task $\displaystyle \frac{\log x}{\log(x-2)} = \frac{\log 9}{\log 3}$
        \task $\displaystyle \frac{\log_2(x-1)}{\log_2(x-3)} = -1$
        \task $\displaystyle \frac{\log_3(2x+1)}{\log_3(x-2)} = 2$
        \task $\displaystyle \frac{2\log(3x-1)}{\log(x+4)} = 1$
        \task $\displaystyle \frac{\log(x^2-5x+6)}{\log(x-1)} = 2$
        \task $\displaystyle \frac{\log_2(x^2-1)}{\log_2(x-1)} = 3$
        \task $\displaystyle \frac{\log(x+1)}{\log(x-2)} = 2$
    \end{tasks}
\end{exercise}


% =========================
% ÚLOHA 14 – Řešte složitější logaritmické rovnice
% =========================
\begin{exercise}[Řešte složitější logaritmické rovnice]
    Řešte v množině reálných čísel $x$:
    \begin{tasks}(2)
        \task $\left(\log_{2} x\right)^{2} + 2\log_{2} x - 3 = 0$
        \task $4\log_{9} x \,(\log_{9} x - 1) = 2 + 3\log_{9} x$
        \task $\left(\log_{\frac12}(x+1)\right)^{2} + 5\log_{\frac12}(x+1) = 6$
        \task $\left(\log_{2}(x-1)\right)^{2} - 3\log_{2}(x-1) + 2 = 0$
        \task $2\left(\log_{3}(x+2)\right)^{2} - \log_{3}(x+2) - 1 = 0$
        \task $2\left(\log_{4}(x-3)\right)^{2} + 3\log_{4}(x-3) - 2 = 0$
        \task $\left(\log_{2}(x+5)\right)^{2} - 4\log_{2}(x+5) + 3 = 0$
        \task $3\left(\log_{7}(x-1)\right)^{2} - 2\log_{7}(x-1) - 1 = 0$
        \task $\left(\log_{10}(x+4)\right)^{2} + \log_{10}(x+4) - 6 = 0$        
        \task \text{$\left(\log_{5}(x^{2}-4)\right)^{2} - 3\log_{5}(x^{2}-4) + 2 = 0$}
        \task $\left(\log_{\frac13}(x-2)\right)^{2} + \log_{\frac13}(x-2) = \log_{\frac13} 9$
        \task $4\log_{3}(2x+1) + \log_{3}\sqrt{2x+1} = \frac{3}{2}\left(\log_{3}(2x+1)\right)^{2}$
    \end{tasks}
\end{exercise}


% =========================
% ÚLOHA 15 – Rovnice s logaritmem v exponentu
% =========================
\begin{exercise}[Řešte rovnice s logaritmem v exponentu]
    Řešte v množině reálných čísel $x$:
    \begin{tasks}(2)
        \task $x^{1+\log x} = 100$
        \task $\sqrt[5]{x^{\log_3 x}} = 243$
        \task $x^{\sqrt{x}} = (\sqrt{x})^{x}$
        \task $(\sqrt{x})^{\log_x (x-2)} = x$
        \task $x^{\log_2 x} = 16$
        \task $x^{\log_5 (x+4)} = 125$
        \task $(x-1)^{\log_x 9} = 27$
        \task $x^{\log_4 (x-3)} = 64$
    \end{tasks}
\end{exercise}


% =========================
% ÚLOHA 16 – Exponenciálně-logaritmické rovnice
% =========================
\begin{exercise}[Řešte exponenciálně-logaritmické rovnice]
    Řešte v množině reálných čísel $x$:
    \begin{tasks}(2)
        \task $x + \log_{2}\!\left(8 + 2^{x}\right) = 7$
        \task $x + \log_{2}\!\left(1 - 3\cdot 2^{x}\right) = x \log_{2} 4$
        \task $\log_{3}\!\left(4\cdot 2^{x}\right) + \log_{3} 3^{x+2} = \log_{3} 36^{x}$
        \task $\log_{5}\!\left(10\cdot 25^{x}\right) - \log_{5}\!\left(5^{x} + 25\right) = x + 1$
        \task $\log_{2}\!\left(5\cdot 2^{x} - 8\right) = x + 2$
        \task $\log_{3}\!\left(2\cdot 3^{x} - 3\right) = x$
        \task $\log_{4}\!\left(2^{x+1} + 8\right) = x$
        \task \text{$\log_{2} 3 + \log_{2} 4^{\,x+\sqrt{x}} = \log_{2}\!\left(2^{x+\sqrt{x}+1} + 4\right) + 2$}
        \task $\log_{5}\!\left(3\cdot 5^{x} - 10\right) + 1 = x$
        \task $\log_{2}\!\left(4 + 2^{x+1}\right) - x = 3$
    \end{tasks}
\end{exercise}


\setcounter{section}{2}
% =========================
% ÚLOHA 17 – Řešte logaritmické nerovnice
% =========================
\begin{exercise}[Řešte následující logaritmické nerovnice]
    Řešte v množině reálných čísel $x$:
    \begin{tasks}(2)
        \task $\log_{3}(x-1) > \log_{3} 3$
        \task $\log_{5}(2x+3) \leq 1$
        \task $\log_{10}(x-3) + \log_{10}(2x-1) < 0$
        \task $\log_{2}(x+1) \geq \log_{2}(5 - x)$
        \task $2\log_{4}(x-1) - \log_{4}(x+3) > 1$
        \task $\log_{0{,}5}(x-2) + \log_{0{,}5}(x-4) \leq -2$
        \task $\log_{3}(x+2) < \log_{3}(7 - x)$
        \task $3\log_{2}(x-1) + \log_{2}(x+3) \geq 5$
        \task $\log_{7}(x^{2} - 4) - \log_{7}(x - 2) < 1$
        \task $2\log_{3}(x + 4) - \log_{3}(x - 1) \leq 2$
        \task $\log_{2}(x - 3) + \log_{2}(x + 5) > 4$
        \task $\log_{5}(2x + 1) - \log_{5}(x - 1) \geq 1$
    \end{tasks}
\end{exercise}


% =========================
% ÚLOHA 18 – Logaritmické nerovnice
% =========================
\begin{exercise}[Řešte logaritmické nerovnice]
    Řešte v množině reálných čísel $x$:
    \begin{tasks}(2)
        \task $\log_{3}(2x - 1) \ge \log_{3}(4x + 3)$
        \task $\log_{\frac{1}{3}}(x + 4) \le \log_{\frac{1}{3}}(5x - 4)$
        \task $\log_{20}(x^{2} - 3x) \le \log_{20} x + \log_{20} 5$
        \task $\log_{0{,}2}(x - 3) > \log_{0{,}2} x - \log_{0{,}2} 2$
        \task $\log_{4}(x - 1) + \log_{4}(x + 3) > 2$
        \task $\log_{5}(x + 2) - \log_{5}(x - 3) < 1$
        \task $\log_{\frac{1}{2}}(x + 1) \ge \log_{\frac{1}{2}}(6 - x)$
        \task $\log_{7}(x - 4) + \log_{7}(x + 2) \le 2$
        \task $\log (2x - 5) - \log (x + 5) > -1$
        \task $2\log_{5}(x - 2) - \log_{5}(x + 1) \le 1$
        \task $\log_{\frac{1}{2}}(x - 1) + \log_{\frac{1}{2}}(x - 5) > -3$
        \task $\log_{3}(2x - 1) - \log_{3}(x + 2) \ge \log_{3} 3$
    \end{tasks}
\end{exercise}


\setcounter{section}{3}
% =========================
% ÚLOHA 19 – BONUS
% =========================
\begin{exercise}[BONUS: Dokažte vztah mezi logaritmy a exponenty]
    Je dáno, že $x = 2^{p}$ a $y = 4^{q}$, kde $p, q \in \mathbb{R}$.
    Dokažte, že platí
    \[
        \log_{2}\left(x^{3} \cdot y\right) = 3p + 2q.
    \]
    Svůj postup přehledně zdůvodněte pomocí vzorců pro logaritmus mocniny
    a vhodného zápisu čísla $4$ jako mocniny čísla $2$.
\end{exercise}


% =========================
% ÚLOHA 20 – BONUS
% =========================
\begin{exercise}[BONUS: Upravte rovnici exponenciální funkce]
    Exponenciální křivka má rovnici
    \[
        y = a b^{x}, \qquad x \in \mathbb{R},
    \]
    kde $a$ a $b$ jsou nenulové konstanty.

    Upravte tuto rovnici tak, aby neznámá $x$ byla vyjádřena jako subjekt
    (tj. vyjádřete ji explicitně v závislosti na ostatních veličinách).
    Konečný tvar výsledku napište pomocí logaritmů o základu $10$.
\end{exercise}


% =========================
% ÚLOHA 22 – BONUS
% =========================
\begin{exercise}[BONUS: Vyjádření neznámé bez logaritmů]
Je dáno, že reálné číslo $x$ splňuje logaritmickou rovnici
\[
    \log_{a} x = 2\bigl(\log_{a} k - \log_{a} 2\bigr),
\]
kde $k>0$, $a>0$ a $a \neq 1$.

\begin{enumerate}[label=\alph*)]
    \item Vyjádřete $x$ pomocí $k$ tak, aby výsledný tvar \emph{neobsahoval logaritmy}.
\end{enumerate}

Nyní předpokládejme, že číslo $x$ splňuje rovnici
\[
    \log_{x} (5y + 1) = 4 + \log_{x} 3,
\]
kde $x>0$, $x \neq 1$ a $y>0$, $y \neq 1$.

\begin{enumerate}[label=\alph*),start=2]
    \item Řešte tuto rovnici pro $y$ a vyjádřete $y$ pomocí $x$ tak, aby výsledný tvar \emph{neobsahoval logaritmy}.
\end{enumerate}
\end{exercise}


% =========================
% ÚLOHA 23 – Exponenciální rovnice (zaokrouhlení na 3 platné číslice)
% =========================
\begin{exercise}[BONUS: Řešte exponenciální rovnice]
    Řešte následující exponenciální rovnice. Výsledky zaokrouhlete na tři platné číslice.
    \begin{tasks}(2)
        \task $5^{2x - 1} = 4^{300}$
        \task $2^{y+1} = \dfrac{10}{2^{y}}$
    \end{tasks}
\end{exercise}


% =========================
% ÚLOHA 24 – BONUS: Úprava složitého logaritmického výrazu
% =========================
\begin{exercise}[BONUS: Zjednodušte logaritmický výraz]
    Zjednodušte co nejvíce následující logaritmický výraz a přitom přehledně ukažte všechny kroky úprav:
    \[
        \log\bigl(10 + 3\sqrt{10}\bigr)
        \;+\;
        \log\bigl(10 + \sqrt{90 + \sqrt{90}}\bigr)
        \;+\;
        \log\bigl(10 - \sqrt{90 + \sqrt{90}}\bigr).
    \]
    Ve svém postupu využijte vzorce pro logaritmus součinu a vhodné úpravy odmocnin.
\end{exercise}


% =========================
% ÚLOHA 25 – BONUS
% =========================
\begin{exercise}[BONUS: Najděte přesnou hodnotu $y$]
    Najděte přesnou hodnotu $y$, která splňuje rovnici
    \[
        2\log\!\left(\frac{x}{y}\right) - 1
        \;=\;
        \log\!\left(10x^{2} \cdot y\right),
        \qquad x \ne 0,\; y \ne 0.
    \]
    Výsledek vyjádřete přesně, bez použití desetinného zápisu.
\end{exercise}


% =========================
% ÚLOHA 26 – BONUS
% =========================
\begin{exercise}[BONUS: Určete parametry mocninné funkce]
    Body $(2,\,10)$ a $(6,\,100)$ leží na křivce dané rovnicí
    \[
        y = a x^{n},
        \qquad a \ne 0,\; n \ne 0.
    \]
    Určete hodnotu $a$ a hodnotu $n$ (zaokrouhlené na tři desetinná místa),
    které tuto křivku určují.
\end{exercise}


% =========================
% ÚLOHA 27 – BONUS
% =========================
\begin{exercise}[BONUS: Rovnice bez reálného řešení]
    Ukažte, že logaritmická rovnice
    \[
        \log_{x^{2} + 2}\!\bigl(2x^{4} - 2x^{3} + 7x^{2} - 2x + 5\bigr) = 2,
        \qquad x \in \mathbb{R},
    \]
    nemá žádné reálné řešení.

    Ve svém postupu nezapomeňte využít podmínky pro základ a argument
    logaritmu a vhodně upravit rovnici do tvaru mocninné rovnice.
\end{exercise}


\setcounter{section}{4}
% =========================
% ÚLOHA 28 – BONUS (\star)
% =========================
\begin{exercise}[$\star$ BONUS: Určete hodnotu logaritmu]
    Jsou dány rovnice
    \[
        \log(xy^{3}) = 1,
        \qquad
        \log(x^{2}y) = 1,
    \]
    kde předpokládáme $x>0$, $y>0$.
    
    Určete hodnotu výrazu
    \[
        \log(xy).
    \]

    Svůj postup přehledně zdůvodněte pomocí vzorců pro součin a mocninu uvnitř logaritmu.
\end{exercise}


% =========================
% ÚLOHA 29 – BONUS (\star)
% =========================
\begin{exercise}[$\star$ BONUS: Určete hodnotu součinu $abcd$]
    Nechť reálná čísla $a, b, c, d$ splňují následující vztahy:
    \[
        4^{a} = 5, \qquad
        5^{b} = 6, \qquad
        6^{c} = 7, \qquad
        7^{d} = 8.
    \]
    Určete hodnotu součinu
    \[
        a \cdot b \cdot c \cdot d.
    \]
    Svůj postup pečlivě zdůvodněte pomocí logaritmů a vlastností exponenciálních funkcí.
\end{exercise}


% =========================
% ÚLOHA 30 – BONUS
% =========================
\begin{exercise}[BONUS: Určete základ logaritmu]
    Řešení rovnice
    \[
        7^{\,x+7} = 8^{\,x}
    \]
    lze zapsat ve tvaru
    \[
        x = \log_{b}\!\left(7^{7}\right),
    \]
    kde $b$ je vhodně zvolený kladný základ logaritmu, $b \neq 1$.

    Určete hodnotu $b$.
\end{exercise}


% =========================
% ÚLOHA 31 – BONUS
% =========================
\begin{exercise}[BONUS: Vyjádření $x^{3}+y^{3}$ pomocí $10^{z}$]
    Nechť $S$ je množina všech uspořádaných trojic $(x,y,z)$ reálných čísel, pro které platí
    \[
        \log_{10}(x + y) = z
        \qquad\text{a}\qquad
        \log_{10}(x^{2} + y^{2}) = z + 1.
    \]

    Existují reálná čísla $a$ a $b$ taková, že pro všechny trojice $(x,y,z) \in S$ platí
    \[
        x^{3} + y^{3} = a \cdot 10^{3z} + b \cdot 10^{2z}.
    \]

    Určete hodnotu součtu
    \[
        a + b.
    \]
\end{exercise}